\section*{Bab \ref{ch:g2:measuring-length} --- Mengukur Panjang}

\begin{solution}{exo:g2:measuring-length:1}
Kalau dekat: berdirilah saling membelakangi. Kalau jauh: ukurlah
dirimu dalam \omterm{def:g2:measuring-length:units}{sentimeter} lalu kirimkan bilangannya --- sebuah bilangan
dapat bepergian di dalam surat, perbandingan punggung ke punggung
tidak.
\end{solution}

\begin{solution}{exo:g2:measuring-length:2}
Tidak ada yang salah hitung. Langkah Mara lebih pendek, sehingga lebih
banyak langkahnya yang termuat di ruang kelas yang sama. \omterm{def:g2:measuring-length:measure}{Satuan} yang
berbeda memberi bilangan yang berbeda untuk panjang yang sama.
\end{solution}

\begin{solution}{exo:g2:measuring-length:3}
\omterm{def:g2:measuring-length:units}{Sentimeter} sama bagi semua orang --- \omterm{def:g2:measuring-length:measure}{satuan} itu bukan milik tubuh
siapa pun. Pengukuran dalam \omterm{def:g2:measuring-length:units}{sentimeter} dapat diperiksa siapa saja
dengan penggaris mana saja.
\end{solution}

\begin{solution}{exo:g2:measuring-length:4}
Semut: \unit{cm} (bahkan kurang!); pintu: \unit{m}; halaman sekolah:
\unit{m}; lebar buku ini: \unit{cm}.
\end{solution}

\begin{solution}{exo:g2:measuring-length:5}
Tidak --- biasanya ada sedikit ruang kosong antara tepi penggaris dan
tanda nolnya, sehingga pensil itu sedikit lebih pendek daripada $13$
\omterm{def:g2:measuring-length:units}{sentimeter}. Tim lupa menempatkan \emph{tanda nol}, bukan tepinya, pada
ujung pensil.
\end{solution}

\begin{solution}{exo:g2:measuring-length:6}
Jawaban beragam; biasanya: lebar tangan kira-kira \qty{7}{cm}, sendok
kira-kira \qty{15}{cm}, cangkir kira-kira \qty{9}{cm}. Yang penting:
setiap bilangan membawa \omterm{def:g2:measuring-length:measure}{satuan} \unit{cm}-nya.
\end{solution}

\begin{solution}{exo:g2:measuring-length:7}
$100 + 100 = 200$: pintu itu tingginya \qty{200}{cm}.
\end{solution}

\begin{solution}{exo:g2:measuring-length:8}
Jawaban beragam; kebanyakan tempat tidur mendekati \qty{2}{m}.
Intinya adalah membandingkan tebakan dengan bilangan hasil ukuran.
\end{solution}

\begin{solution}{exo:g2:measuring-length:9}
Pita \qty{45}{cm} yang lebih panjang: $45 - 38 = 7$, jadi selisihnya
\qty{7}{cm}. Tanpa bilangan, menjajarkan keduanya juga menemukan yang
lebih panjang --- tetapi ``\qty{45}{cm}'' muat dalam sebuah surat,
sedangkan ``dijajarkan'' tidak.
\end{solution}

\begin{solution}{exo:g2:measuring-length:10}
Tidak ada orang lain yang punya kaki Kakek: anak kecil yang melangkahi
kebun yang sama mungkin menghitung $60$ langkah kecil, dan seorang
tamu tidak dapat memeriksa ``$30$ langkah'' tanpa kehadiran Kakek.
Diukur dalam \omterm{def:g2:measuring-length:units}{meter} --- yang sama bagi semua orang --- panjang kebun
itu dapat dicatat, dikirim, dan diperiksa siapa pun.
\end{solution}
